\lyseis{Β1.8}{Παραπληρωματικές, συμπληρωματικές, κατακορυφήν}
\sol{1} $55^\circ$, $28^\circ$, $82^\circ$, $45^\circ30'$.
\sol{2} $50^\circ$, $90^\circ$, $175^\circ$, $106^\circ45'$.
\sol{3} $48^\circ$ (κατακορυφήν), $132^\circ$ και $132^\circ$.
\sol{4} Όλες ορθές.
\sol{5} Η $45^\circ$ ($45+45=90$). Η $90^\circ$ ($90+90=180$).
\sol{6} $72^\circ$ και $108^\circ$.
\sol{7} α) Σωστό. β) Λάθος: παραπληρωματικές. γ) Λάθος. δ) Σωστό. ε) Σωστό.
\sol{8} $70^\circ$, $110^\circ$, $70^\circ$, $110^\circ$. Δύο διαφορετικές τιμές.
\sol{9} $90^\circ-35^\circ=55^\circ$.
\sol{10} $2x+10=3x-20$, $x=30$ και $\omega=70^\circ$.
\sol{11} $90^\circ-72^\circ=18^\circ$.
\sol{12} Κατακορυφήν γωνίες ίσες: $40^\circ$.
\sol{13} Η $\omega$ και η διπλανή της $\varphi$ έχουν $\omega+\varphi=180^\circ$. Η κατακορυφήν της $\omega$ (έστω $\omega'$) έχει επίσης $\omega'+\varphi=180^\circ$. Άρα $\omega=\omega'$.
\sol{14} Όχι. Δύο γωνίες $50^\circ$ σε διαφορετικά σημεία είναι ίσες αλλά δεν είναι κατακορυφήν.
\sol{15} Η παραπληρωματική της οξείας είναι αμβλεία. Η συμπληρωματική είναι οξεία.
\sol{16} Σχηματίζουν $90^\circ$: $25^\circ+65^\circ=90^\circ$ (οι μισές είναι $25^\circ$ και $65^\circ$).
\sol{17} Ελεύθερη απάντηση.
\sol{18} Ελεύθερη απάντηση.
